% Build with: latexmk -pdf main.tex
\documentclass[11pt]{article}

\usepackage[T1]{fontenc}
\usepackage{lmodern}
\usepackage[margin=1in]{geometry}
\usepackage{microtype}
\usepackage{amsmath,amssymb}
\usepackage{booktabs}
\usepackage{array}
\usepackage{enumitem}
\usepackage{xcolor}
\usepackage[colorlinks=true,linkcolor=blue!50!black,citecolor=blue!50!black,urlcolor=blue!50!black]{hyperref}

\newcommand{\bits}[1]{\{0,1\}^{#1}}
\newcommand{\getsr}{\stackrel{\$}{\gets}}
\newcommand{\keccak}{\mathsf{Keccak}}
\newcommand{\Th}{\mathsf{Th}}
\newcommand{\PRF}{\mathsf{PRF}}
\newcommand{\PRFmsg}{\mathsf{PRF}_{\mathrm{msg}}}
\newcommand{\Hmsg}{\mathsf{H}_{\mathrm{msg}}}
\newcommand{\Enc}{\mathsf{Enc}}
\newcommand{\Gen}{\mathsf{Gen}}
\newcommand{\Sig}{\mathsf{Sig}}
\newcommand{\Ver}{\mathsf{Ver}}
\newcommand{\BE}{\mathsf{BE}}
\newcommand{\Truncate}{\mathsf{Trunc}}
\newcommand{\basetwo}{\mathsf{base\_2^b}}
\newcommand{\concat}{\mathbin\Vert}
\newcommand{\PK}{\mathit{PK}}
\newcommand{\SK}{\mathit{SK}}
\newcommand{\pkseed}{\mathit{PK.seed}}
\newcommand{\pkroot}{\mathit{PK.root}}
\newcommand{\skseed}{\mathit{SK.seed}}
\newcommand{\skprf}{\mathit{SK.prf}}
\newcommand{\addr}{\mathit{ADRS}}
\newcommand{\cmax}{C_{\max}}
\newcommand{\idx}{\mathit{idx}}
\newcommand{\WotsSign}{\mathsf{Wots.sign}}
\newcommand{\WotsLeaf}{\mathsf{Wots.leaf}}
\newcommand{\TreeFold}{\mathsf{Tree.fold}}
\newcommand{\ForsKey}{\mathsf{Fors.key}}
\newcommand{\ForsSign}{\mathsf{Fors.sign}}
\newcommand{\ForsRec}{\mathsf{Fors.recover}}

\emergencystretch=1.5em
\setlist[enumerate]{leftmargin=2em, itemsep=4pt}

\title{leanSPHINCS}
\author{}
\date{}

\begin{document}
\maketitle

\begin{abstract}

This specification defines leanSPHINCS, a stateless hash-based signature scheme: SLH-DSA~\cite{FIPS205} on Keccak-256, with WOTS$^+$C~\cite{HK22C} as its one-time signature and a hypertree of three layers of different heights. It keeps FIPS~205's algorithms, address and input layouts wherever the table of Section~\ref{sec:overview} does not say otherwise, but it is not a FIPS~205 parameter set and does not interoperate with FIPS~205 implementations.

\begin{itemize}
  \item \textbf{NIST security level~1}~\cite{NISTPQC}: a target of 127 bits of classical security in the random-oracle model and about 64 bits in the quantum random-oracle model (Section~\ref{sec:security}).
  \item \textbf{Public key: 32 bytes.}
  \item \textbf{Signature: 7740 bytes.}
  \item \textbf{Verification: 515 permutations} of Keccak-f.
  \item \textbf{Lifetime: $2^{30}$ signatures per key pair.}
\end{itemize}
\end{abstract}

\section{Overview and differences from FIPS 205}
\label{sec:overview}

A message digest picks a FORS instance and one leaf in each of its trees. The FORS signature signs the message, and a hypertree of WOTS$^+$C keys signs the FORS public key up to the public key. The table below lists every point where leanSPHINCS departs from FIPS~205. Everything else is FIPS~205's: the address and its types $0$ to $6$; the input layouts of $\mathsf F$, $\mathsf H$, $\mathsf T_\ell$, $\PRF$ and $\PRFmsg$; FORS, its secrets, its trees and the compression of its roots by $\mathsf T_k$; the hashing of the hypertree's nodes; the split of the digest into FORS leaves, tree index and leaf index; pure SLH-DSA's framing of the message with its context; deterministic signing with $\mathit{opt\_rand}=\pkseed$; the public key $\pkseed\concat\pkroot$; and the order of a signature.

\begin{table}[ht]
\centering
\small
\renewcommand{\arraystretch}{1.15}
\begin{tabular}{@{}>{\raggedright}p{0.15\linewidth}>{\raggedright}p{0.33\linewidth}>{\raggedright\arraybackslash}p{0.46\linewidth}@{}}
\toprule
& FIPS~205 & leanSPHINCS\\
\midrule
Hash & SHAKE256 with output of $8n$ or $8m$ bits & Keccak-256, the first $n$ bytes of its output (Section~\ref{sec:parameters})\\
$\Hmsg$ & one SHAKE256 call of $m$ bytes & $\mathsf{MGF1}$~\cite{RFC8017} on Keccak-256: two calls with a 32-bit counter appended, $m$ bytes of their outputs\\
One-time signature & WOTS$^+$: $\mathit{lg}_w=4$, $2n+3$ chains with a checksum & WOTS$^+$C: $w=2$, $v=64$ chains whose positions sum to $T$, a counter per layer searched by the signer, an encoding hash with address type $7$ (Section~\ref{sec:ots})\\
Hypertree & $d$ layers of height $h'=h/d$ & three layers of heights $8$, $8$, $10$ from the bottom; the leaf index has $h_0$ bits (Section~\ref{sec:layer})\\
Parameters & the approved sets & $n=16$, $h=26$, $d=3$, $k=24$, $a=10$, $m=34$\\
Keys & $\skseed$, $\skprf$, $\pkseed$ drawn independently, $n$ bytes each; the secret key holds all four values & $\skseed$, 32 bytes, is the whole secret key; $\pkseed\concat\skprf$ is its hash (Section~\ref{sec:keys}); $\PRF$'s input is 80 bytes\\
Signing & hedged or deterministic; any message, any context, or a pre-hash & deterministic pure SLH-DSA only, on a 32-byte message under the fixed context \texttt{leanSPHINCS}\\
Signature & $R$, $\mathit{SIG}_{\mathrm{FORS}}$, $\mathit{SIG}_{\mathrm{HT}}$ & the same, each layer's counter $\BE_{32}(c_l)$ before its chain values\\
Lifetime & $2^{64}$ signatures & $2^{30}$ signatures\\
\bottomrule
\end{tabular}
\end{table}

\section{Hashing and parameters}
\label{sec:parameters}

Write $\bits r$ for $r$-bit strings, $\concat$ for concatenation, $\BE_r(z)$ for the unsigned $r$-bit big-endian encoding of $z$ (FIPS~205's $\mathsf{toByte}(z,r/8)$), $\BE^{-1}$ for its inverse ($\mathsf{toInt}$), and $x\getsr A$ for a uniform sample from $A$. Indices start at zero; $\bot$ denotes failure.

Let $\keccak:\bits{*}\to\bits{256}$ be Keccak-256~\cite{keccak}, Ethereum's hash: the 256-bit hash of FIPS~202~\cite{fips202} but for its padding's first byte, $\mathtt{0x01}$ where the standard has $\mathtt{0x06}$. It takes the place of SHAKE256, and $\Truncate_r$ keeps the first $r$ bytes of its output. With $\addr$ a 32-byte address (Appendix~\ref{sec:tweaks}) naming one hash call in the whole structure,
\[
\begin{aligned}
  \Th(\pkseed,\addr,M)&=\Truncate_n\!\left(\keccak(\pkseed\concat\addr\concat M)\right),\\
  \PRF(\pkseed,\skseed,\addr)&=\Th(\pkseed,\addr,\skseed),\\
  \PRFmsg(\skprf,\mathit{opt\_rand},M)&=\Truncate_n\!\left(\keccak(\skprf\concat\mathit{opt\_rand}\concat M)\right),
\end{aligned}
\]
$\Th$ being FIPS~205's $\mathsf F$, $\mathsf H$ and $\mathsf T_\ell$ alike, which the text below names as FIPS~205 does. As pure SLH-DSA does (FIPS~205 Algorithms~22 and~24), $\PRFmsg$ and $\Hmsg$ take the message framed with a context string $\mathit{ctx}$, here the 11 ASCII bytes of \texttt{leanSPHINCS}, so that $M'$ is 45 bytes:
\[
  M'=\BE_8(0)\concat\BE_8(|\mathit{ctx}|)\concat\mathit{ctx}\concat M.
\]
The message hash needs more output than one call gives, so it is $\mathsf{MGF1}$~\cite{RFC8017} on Keccak-256: two calls on $Z=R\concat\pkseed\concat\pkroot\concat M'$ and a counter, of 97 bytes each,
\[
  \Hmsg(R,\pkseed,\pkroot,M')=\Truncate_m\!\left(\keccak(Z\concat\BE_{32}(0))\concat\keccak(Z\concat\BE_{32}(1))\right).
\]

The message $M$ is 32 bytes. $\pkseed$, $\pkroot$, $\skprf$, the randomizer $R$, chain values and Merkle nodes are $n$ bytes; $\skseed$ is 32 bytes.

\begin{center}
\begin{tabular}{@{}lll@{}}
\toprule
Symbol & Value & Meaning\\
\midrule
$n$ & $16$ & bytes of a hash value\\
$w$ & $2$ & bits of a chain position (FIPS~205's $\mathit{lg}_w$)\\
$v$ & $64$ & chains of a one-time key, $8n/w$\\
$T$ & $128$ & sum of the signed chain positions\\
$d$ & $3$ & layers of the hypertree\\
$h_0,h_1,h_2$ & $8,8,10$ & height of a tree of each layer, from the bottom\\
$h$ & $26$ & $h_0+h_1+h_2$\\
$k$ & $24$ & trees of a FORS instance\\
$a$ & $10$ & height of a FORS tree\\
$m$ & $34$ & bytes of $\Hmsg$: $\lceil ka/8\rceil+\lceil(h-h_0)/8\rceil+\lceil h_0/8\rceil$\\
$\cmax$ & $2^{32}$ & encoding attempts per layer\\
\bottomrule
\end{tabular}
\end{center}

\paragraph{Keys.}
\label{sec:keys}
The secret key is $\skseed$ alone, 32 bytes, so that guessing it costs $2^{256}$ hashes. The other two seeds are its hash under the scheme's 11-byte ASCII name, so that no other scheme derives from the same bytes:
\[
  \pkseed\concat\skprf=\keccak\!\left(\text{\texttt{leanSPHINCS}}\concat\skseed\right).
\]
Functions that read secrets use $\skseed$, and all hashes $\pkseed$, implicitly.

\section{\texorpdfstring{WOTS$^+$C}{WOTS+C}: signing with hash chains}
\label{sec:ots}

A one-time key $K=(l,\tau,e)$ is identified by its layer $l<d$, counted from the bottom, its tree $\tau$ on that layer and its leaf $e<2^{h_l}$ in that tree, FIPS~205's key pair address. It has $v$ chains of $2^w=4$ positions, $0$ through $3$, computed as FIPS~205's WOTS$^+$ computes them: for $0\leq i<v$,
\[
\begin{aligned}
  C_{i,0}&=\PRF\!\left(\pkseed,\skseed,\addr_{\mathrm{prf}}(K,i)\right),\\
  C_{i,\mu+1}&=\mathsf F\!\left(\pkseed,\addr_{\mathrm{chain}}(K,i,\mu),C_{i,\mu}\right),
    &&0\leq\mu<2^w-1,
\end{aligned}
\]
a step's address holding the position it leaves. The key's public key, its Merkle leaf, compresses the endpoints in chain order:
\[
  X^{l,\tau}_{0,e}=\mathsf T_v\!\left(\pkseed,\addr_{\mathrm{pk}}(K),C_{0,3}\concat\cdots\concat C_{v-1,3}\right).
\]

To sign an $n$-byte $M$, a FORS public key (Section~\ref{sec:fts}) or the root of a tree of the layer below, the signer encodes it as positions $x_0,\ldots,x_{v-1}$ and reveals $C_{i,x_i}$ from each chain; the verifier hashes these forward to the endpoints. In place of WOTS$^+$'s checksum, every valid encoding has the same sum:
\[
  \mathcal C=\left\{x\in\{0,\ldots,2^w-1\}^{v}:\textstyle\sum_{i}x_i=T\right\}.
\]
Hashing forward only increases positions, and two distinct encodings of equal sum are incomparable, so the values revealed for one encoding never give those of another.

\paragraph{Encoding.}
$\Enc(K,M,c)$ computes, from an input of 68 bytes,
\[
  D=\Th\!\left(\pkseed,\addr_{\mathrm{enc}}(K),M\concat\BE_{32}(c)\right),
\]
and reads it as FIPS~205's WOTS$^+$ reads a message, $x=\basetwo(D,w,v)$: with $D_j$ byte $j$ of $D$, each byte's most significant bits first,
\[
  x_i=\left\lfloor D_{\lfloor i/4\rfloor}/2^{6-2(i\bmod4)}\right\rfloor\bmod 4,\qquad 0\leq i<v.
\]
It returns $x$ if $x\in\mathcal C$, and $\bot$ otherwise. Since $wv=8n$, an encoding reads every bit of $D$.

\paragraph{Signing.}
$\WotsSign(K,M)$ tries $c=0,\ldots,\cmax-1$ in order and, at the first $c$ with $x=\Enc(K,M,c)\neq\bot$, returns $(c,\sigma)$ with $\sigma_i=C_{i,x_i}$. It returns $\bot$ if every counter fails. Taking the least such counter makes signing one $M$ under one key twice reveal the same positions: different positions for the same $M$ would let an adversary combine the smaller positions of both signatures.

\paragraph{Recovering the leaf.}
$\WotsLeaf(K,M,c,\sigma)$ computes $x=\Enc(K,M,c)$ and returns $\bot$ if it fails. Otherwise it sets $C'_{i,x_i}=\sigma_i$, applies the chain recurrence up to position $3$, and returns $\mathsf T_v\!\left(\pkseed,\addr_{\mathrm{pk}}(K),C'_{0,3}\concat\cdots\concat C'_{v-1,3}\right)$.

\section{The hypertree: authenticating one-time keys}
\label{sec:layer}

Layer $l$ has $2^{h-(h_0+\cdots+h_l)}$ trees of height $h_l$: $2^{18}$ on the bottom layer, $2^{10}$ on the middle one, and a single top tree. The leaves of tree $\tau$ of layer $l$ are the $X^{l,\tau}_{0,e}$, and its inner nodes are those of FIPS~205's $\mathsf{xmss\_node}$,
\[
  X^{l,\tau}_{z,j}=\mathsf H\!\left(\pkseed,\addr_{\mathrm{node}}(l,\tau,z,j),X^{l,\tau}_{z-1,2j}\concat X^{l,\tau}_{z-1,2j+1}\right),\qquad 1\leq z\leq h_l,\ 0\leq j<2^{h_l-z}.
\]
The tree's root is $X^{l,\tau}_{h_l,0}$, and the root of the top tree is $\pkroot$.

An index $\idx<2^h$ names one key on every layer: with $\beta_l=h_0+\cdots+h_{l-1}$ the index bits below layer $l$, its key there is
\[
  K_l(\idx)=\left(l,\ \left\lfloor\idx/2^{\beta_l+h_l}\right\rfloor,\ \left\lfloor\idx/2^{\beta_l}\right\rfloor\bmod2^{h_l}\right).
\]
This is FIPS~205's $\mathsf{ht\_sign}$ with each layer's own height for $h'$: the bottom key is at leaf $\idx\bmod2^8$ of tree $\lfloor\idx/2^8\rfloor$, the middle key at leaf $\lfloor\idx/2^8\rfloor\bmod2^8$ of tree $\lfloor\idx/2^{16}\rfloor$, and the top key at leaf $\lfloor\idx/2^{16}\rfloor$. The key of layer $l+1$ signs the root of the tree holding the key of layer $l$.

The authentication path of the key at leaf $e$ is $P=(P_0,\ldots,P_{h_l-1})$, from the leaf up, with $P_z=X^{l,\tau}_{z,\lfloor e/2^z\rfloor\oplus1}$ the sibling at height $z$, where $\oplus1$ flips the low bit of the node index.

$\TreeFold(K,Q_0,P)$ recovers the root from a leaf value $Q_0$ and its path: for $0\leq z<h_l$, with $j=\lfloor e/2^{z+1}\rfloor$,
\[
  Q_{z+1}=\begin{cases}
    \mathsf H\!\left(\pkseed,\addr_{\mathrm{node}}(l,\tau,z+1,j),Q_z\concat P_z\right)
      &\text{if bit $z$ of $e$ is zero},\\[2pt]
    \mathsf H\!\left(\pkseed,\addr_{\mathrm{node}}(l,\tau,z+1,j),P_z\concat Q_z\right)
      &\text{otherwise},
  \end{cases}
\]
and it returns $Q_{h_l}$.

\section{FORS: signing the digest}
\label{sec:fts}

The FORS instance of an index is FIPS~205's at the bottom key $K_0(\idx)=(0,\tau,e)$: its addresses hold layer $0$, the tree $\tau$ and the key pair address $e$. It has $k$ trees of height $a$, whose leaves hash secrets. A node's tree index counts across the trees: in tree $\kappa<k$, the node at height $z$ and index $j<2^{a-z}$ has tree index $\kappa2^{a-z}+j$. For leaf $s<2^a$ and $1\leq z\leq a$:
\[
\begin{aligned}
  \mathit{sk}_{\kappa,s}&=\PRF\!\left(\pkseed,\skseed,\addr_{\mathrm{fprf}}(\idx,\kappa2^a+s)\right),\\
  Y^{\kappa}_{0,s}&=\mathsf F\!\left(\pkseed,\addr_{\mathrm{fors}}(\idx,0,\kappa2^a+s),\mathit{sk}_{\kappa,s}\right),\\
  Y^{\kappa}_{z,j}&=\mathsf H\!\left(\pkseed,\addr_{\mathrm{fors}}(\idx,z,\kappa2^{a-z}+j),Y^{\kappa}_{z-1,2j}\concat Y^{\kappa}_{z-1,2j+1}\right).
\end{aligned}
\]
The instance's public key compresses the trees' roots:
\[
  \ForsKey(\idx)=\mathsf T_k\!\left(\pkseed,\addr_{\mathrm{roots}}(\idx),Y^{0}_{a,0}\concat\cdots\concat Y^{k-1}_{a,0}\right).
\]

\paragraph{Opening.}
The digest selects a leaf $s_\kappa$ in each tree. $\ForsSign(\idx,s)$ reveals, per tree, the secret $\mathit{sk}_{\kappa,s_\kappa}$ and the path $(Y^{\kappa}_{z,\lfloor s_\kappa/2^z\rfloor\oplus1})_{0\leq z<a}$, from the leaf up: $1+a$ values per tree.

\paragraph{Recovering the key.}
$\ForsRec(\idx,s,\text{opening})$, FIPS~205's $\mathsf{fors\_pkFromSig}$, hashes each revealed secret into its leaf, folds it to its tree's root as $\TreeFold$ does, with the bits of $s_\kappa$ and the tree's node addresses, and returns $\mathsf T_k$ of the $k$ roots as $\ForsKey$ does.

An instance may sign several digests: without breaking the hash, a forgery at it needs, in every tree, a selected leaf whose secret is already revealed.

\paragraph{The digest.}
For message $M$ and randomizer $R$, the digest $\Hmsg(R,\pkseed,\pkroot,M')$ has $m=34$ bytes $g_0,\ldots,g_{33}$, which FIPS~205's Algorithm~19 splits into its first $\lceil ka/8\rceil=30$ bytes $\mathit{md}$, the next $\lceil(h-h_0)/8\rceil=3$ bytes, and the last $\lceil h_0/8\rceil=1$ byte. The FORS leaves are $\mathit{md}$ read $a$ bits at a time, most significant first, and the index is the tree, from the middle part, then the leaf, from the last byte:
\[
\begin{aligned}
  (s_0,\ldots,s_{k-1})&=\basetwo(\mathit{md},a,k),\\
  \idx&=2^{h_0}\left(\BE^{-1}(g_{30}\concat g_{31}\concat g_{32})\bmod2^{h-h_0}\right)+\left(g_{33}\bmod2^{h_0}\right).
\end{aligned}
\]
Every digest is accepted, since every index names a bottom leaf.

\section{The signature scheme}
\label{sec:scheme}

\subsection{Key generation}

$\Gen$ samples $\skseed\getsr\bits{256}$, derives $\pkseed$ and $\skprf$ (Section~\ref{sec:keys}), builds the top tree, as FIPS~205's $\mathsf{slh\_keygen\_internal}$ does, and returns
\[
  \PK=(\pkseed,\pkroot),\qquad \SK=\skseed.
\]

\subsection{Signing}

$\Sig(\SK,M)$, on a 32-byte message $M$, is the deterministic $\mathsf{slh\_sign}$, which runs $\mathsf{slh\_sign\_internal}$ on $M'$:
\begin{enumerate}
\item Derive $\pkseed$ and $\skprf$ from $\SK$, and $\pkroot$ from the top tree.
\item Compute the randomizer $R=\PRFmsg(\skprf,\pkseed,M')$, and extract $\idx$ and $s$ from $\Hmsg(R,\pkseed,\pkroot,M')$.
\item Compute the opening $\ForsSign(\idx,s)$ and $M_0=\ForsKey(\idx)$.
\item For $l=0,1,2$: compute $(c_l,\sigma_l)=\WotsSign(K_l(\idx),M_l)$, returning $\bot$ on failure, and the path $P_l$ of $K_l(\idx)$ in its tree, and let $M_{l+1}$ be that tree's root.
\item Return $R$, the opening, and $(c_l,\sigma_l,P_l)$ for each layer.
\end{enumerate}
Signing changes no state, and repeating a message returns the same signature, or the same failure. A key signs at most $2^{30}$ messages.

\subsection{Verification}

$\Ver(\PK,M,\sigma)$ rejects a signature of the wrong length. Otherwise it parses $\PK=(\pkseed,\pkroot)$ and $\sigma$ (Section~\ref{sec:costs}):
\begin{enumerate}
\item Extract $\idx$ and $s$ from $\Hmsg(R,\pkseed,\pkroot,M')$.
\item Set $M_0=\ForsRec(\idx,s,\text{opening})$.
\item For $l=0,1,2$: compute $X_l=\WotsLeaf(K_l(\idx),M_l,c_l,\sigma_l)$, rejecting if $X_l=\bot$, and $M_{l+1}=\TreeFold(K_l(\idx),X_l,P_l)$.
\item Accept exactly when $M_3=\pkroot$.
\end{enumerate}

\subsection{Serialization and cost}
\label{sec:costs}

The public key is $\pkseed\concat\pkroot$, 32 bytes, and the secret key $\skseed$, 32 bytes. A signature is $R$, then the FORS opening tree by tree, each tree's secret then its path from the leaf up, then the layers from the bottom up, each $\BE_{32}(c_l)$, then its $v$ chain values in chain order, then its path from the leaf up:
\[
  \underbrace{n}_{R}
  +\underbrace{k\,(1+a)\,n}_{\text{FORS}}
  +\underbrace{d\,(4+v\,n)}_{\text{WOTS$^+$C}}
  +\underbrace{h\,n}_{\text{paths}}
  =16+4224+3084+416=7740\text{ bytes}.
\]

Costs count permutations of Keccak-f: an input of $r$ bytes costs $\lfloor r/136\rfloor+1$. An addressed input starts with the 48 bytes of $\pkseed\concat\addr$. Thus $\mathsf F$ (64 bytes), $\mathsf H$ and $\PRF$ (80), an encoding (68), $\PRFmsg$ (77) and each call of $\Hmsg$ (97) cost one permutation, $\mathsf T_k$ (432) four, and $\mathsf T_v$ (1072) eight. A verification walks exactly $(2^w-1)v-T=64$ chain steps per layer, so it takes exactly
\[
  \underbrace{2}_{\Hmsg}
  +\underbrace{k\,(1+a)}_{\text{FORS trees}}
  +\underbrace{4}_{\mathsf T_k}
  +\underbrace{d}_{\text{encodings}}
  +\underbrace{d\bigl((2^w-1)v-T+8\bigr)}_{\text{WOTS$^+$C}}
  +\underbrace{h}_{\text{paths}}
  =2+264+4+3+216+26=515
\]
permutations.

A one-time leaf takes $v$ calls of $\PRF$, $(2^w-1)v=192$ chain steps and $\mathsf T_v$: 264 permutations. Key generation takes $264\cdot2^{10}+(2^{10}-1)+1=271\,360$, the last for the derivation of the seeds. Signing takes, in expectation,
\[
  \underbrace{4}_{\text{seeds, }\PRFmsg,\,\Hmsg}
  +\underbrace{271\,359}_{\text{top tree}}
  +\underbrace{73\,708}_{\text{FORS}}
  +\underbrace{135\,678}_{\text{lower trees}}
  +\underbrace{d\,(v+T+1/p)}_{\text{WOTS$^+$C}}
  \approx5.25\cdot10^5,
\]
where a one-time signature takes $v$ calls of $\PRF$, $T$ chain steps and $1/p\approx14\,670.4$ expected encoding attempts of one permutation each, with $p$ of Section~\ref{sec:completeness}. A signer may keep the top tree from key generation instead of rebuilding it.

\section{Security}
\label{sec:security}

Model $\keccak$ as one shared random oracle. The challenger generates a key pair and gives $\PK$ to a classical probabilistic adversary, which may adaptively query the oracle and request signatures on messages of its choice, at most $q_s$ of them. It finally returns $(M^*,\sigma^*)$ and wins exactly when $\Ver(\PK,M^*,\sigma^*)$ accepts and no signing request for $M^*$ returned $\sigma^*$: \emph{strong unforgeability}. A budget $q$ bounds every oracle call of the experiment, including key generation, signing and the final verification. The target is that for $q_s=2^{30}$ every adversary wins with probability at most $q/2^{127}$.

\paragraph{The Lean proof.}
The Lean~4 theorem \texttt{FORS.h26\_layered} in \texttt{formal/lean\_sphincs/} proves this bound, for $q_s=1\,182\,000\,000\approx2^{30.14}$, for an earlier version of the scheme with the same parameters and different hash inputs: 16-byte little-endian addresses; a public parameter, every secret and the randomizer derived from a 32-byte hash of the seed, two secrets per call; a digest of two calls over the parameter, an address, the message and the randomizer, read least significant bit first; and FORS trees without root hashes, whose key is their $2k$ top nodes, absorbed whole by the bottom encoding. Its model fixes every hash input's bytes, so it does not cover the scheme of this document.

\paragraph{Hash domains.}
Every call of the hypertree and of FORS hashes $\pkseed$ then an address naming it: its type, among FIPS~205's or the encoding's own, and its fields, so distinct calls of one key never share an input, and calls of different keys differ in $\pkseed$. The calls without an address are $\PRFmsg$, the calls of $\Hmsg$ and the seed derivation, whose 77-byte, 97-byte and 43-byte inputs no other call has the length of.

\paragraph{Quantum security.}
About 64 bits of security in the quantum random-oracle model is a target. No quantum bound is proved.

\section{Completeness}
\label{sec:completeness}

Signing fails only if a counter search is exhausted; whenever it succeeds, verification accepts, each recovery retracing what the signer computed. An encoding accepts exactly when its 64 two-bit positions sum to 128, so a uniform 128-bit $D$ is accepted with probability
\[
  p=\frac{[y^{128}](1+y+y^2+y^3)^{64}}{2^{128}}
  =\frac{23195210617839120199658819909693679}{2^{128}}
  \geq2^{-13.85},
\]
where $[y^r]$ denotes the coefficient of $y^r$, and all $\cmax=2^{32}$ counters of a layer fail with probability $(1-p)^{\cmax}\leq\exp(-2^{18})$.

\appendix
\section{Address encoding}
\label{sec:tweaks}

An address is FIPS~205's (Section~4.2 there), 32 bytes in eight big-endian words: the layer address, the tree address in three words, the type, then three words whose meaning the type gives,
\[
  \mathsf{enc}(l,\tau,t,u_1,u_2,u_3)=\BE_{32}(l)\concat\BE_{96}(\tau)\concat\BE_{32}(t)\concat\BE_{32}(u_1)\concat\BE_{32}(u_2)\concat\BE_{32}(u_3).
\]
For a one-time key $K=(l,\tau,e)$, the key pair address $u_1$ is its leaf $e$, and a FORS instance's addresses are those of its bottom key $K_0(\idx)=(0,\tau,e)$. Types $0$ to $6$ are FIPS~205's; type $7$ is the WOTS$^+$C encoding's.

\begin{center}
\small
\begin{tabular}{@{}lll@{}}
\toprule
Address & FIPS~205 type & Encoding\\
\midrule
$\addr_{\mathrm{chain}}(K,i,\mu)$ & \texttt{WOTS\_HASH} & $\mathsf{enc}(l,\tau,0,e,i,\mu)$\\
$\addr_{\mathrm{pk}}(K)$ & \texttt{WOTS\_PK} & $\mathsf{enc}(l,\tau,1,e,0,0)$\\
$\addr_{\mathrm{node}}(l,\tau,z,j)$ & \texttt{TREE} & $\mathsf{enc}(l,\tau,2,0,z,j)$\\
$\addr_{\mathrm{fors}}(\idx,z,j)$ & \texttt{FORS\_TREE} & $\mathsf{enc}(0,\tau,3,e,z,j)$\\
$\addr_{\mathrm{roots}}(\idx)$ & \texttt{FORS\_ROOTS} & $\mathsf{enc}(0,\tau,4,e,0,0)$\\
$\addr_{\mathrm{prf}}(K,i)$ & \texttt{WOTS\_PRF} & $\mathsf{enc}(l,\tau,5,e,i,0)$\\
$\addr_{\mathrm{fprf}}(\idx,j)$ & \texttt{FORS\_PRF} & $\mathsf{enc}(0,\tau,6,e,0,j)$\\
$\addr_{\mathrm{enc}}(K)$ & none & $\mathsf{enc}(l,\tau,7,e,0,0)$\\
\bottomrule
\end{tabular}
\end{center}

\bibliographystyle{alphaurl}
\bibliography{refs}

\end{document}
